\section{Conclusion}

This study constructed a consolidated financial database consisting of four financial assets (Apple Inc (AAPL), BDO Unibank (BDO), Bitcoin (BTC), and the SPDR S\&P 500 ETF Trust (SPY)) and four macroeconomic indicators from the Federal Reserve Economic Data (FRED) database. The study developed a consistent analytical framework over the 2020--2025 period by systematically acquiring, cleaning, standardizing the calendar, filling in missing data, integrating the data into a database, and conducting analysis using SQL.

The empirical results show that returns are not the only measure for assessing investment performance. Bitcoin had the largest downside risk along with the highest volatility and the highest cumulative return, reflecting the fundamental risk-return trade-off. BDO generated the highest risk-adjusted return and consistently had the lowest correlation with the other assets, highlighting its diversification value. Apple demonstrated strong long-term growth prospects with relatively low risk, while SPY provided diversified exposure to the broader U.S. equity market and served as a stable benchmark for asset performance evaluation. Together, these findings suggest that no single asset dominates in every respect and that a diversified portfolio is likely the best approach to achieving investment goals.

The analyses also demonstrated that asset performance varies under different macroeconomic conditions, such as high market volatility, inflation, and interest-rate fluctuations. The exploratory analyses and SQL investigations repeatedly show that asset performance is influenced differently across macroeconomic regimes, indicating that investors should consider not only past performance but also how assets behave in the wider economic context. Similarly, the modest correlations between BDO and the other assets suggest that diversification continues to play a role in decreasing overall portfolio risk.

Overall, the results confirm the tenets of modern portfolio theory and show that the key to a sound investment portfolio is not necessarily to maximize the performance of any individual asset but to balance expected return, risk, and diversification. Bitcoin exhibited the highest capital appreciation, but its high volatility makes it unsuitable as a primary portfolio holding. Conversely, BDO's comparatively small cumulative return does not diminish its value, since its superior risk-adjusted performance and diversification benefits increase the overall efficiency of the portfolio. Hence, a multi-asset approach with exposure to Apple, BDO, SPY, and a measured allocation to Bitcoin offers a more balanced approach to long-term investing compared to a concentrated focus in a single asset class. Although market outcomes cannot be predicted from historical performance alone, the analytical framework developed in this study demonstrates how integrated financial databases and SQL-based financial analysis can contribute to more informed and evidence-based investment decision making.
